\section{Integration with the BX3 Framework}
\label{sec:rs-integration}

Recursive Spawning does not operate as a standalone delegation mechanism.
Its correctness guarantees are derived entirely from the structural properties
enforced by the BX3 three-layer architecture operating in mandatory concert.
This section specifies the three integration points and establishes the
unified invariant that downstream correctness proofs depend upon.

% ---------------------------------------------------------------------------
\subsection{Purpose Layer as Accountability Anchor}
\label{sec:rs-purpose-anchor}

Every node $n$ in a BX3 spawn tree carries a human anchor $h(n)$ established
at node provisioning and structurally non-modifiable for the node's entire
operational lifetime.  The \Purpose{} layer is the architectural locus of this
anchor: it encodes the principal's intent, defines the operating range $R(n)$,
and serves as the final machine layer before the human principal in any
escalation chain.

Three integration points connect Recursive Spawning to the \Purpose{} layer:

\begin{enumerate}
  \item[(a)] \textbf{Anchor inheritance at spawn.}  When parent node $p$ spawns
    child $c$, the spawn directive encodes $h(c) = h(p)$.  The \Purpose{} layer
    of $c$ inherits the same human anchor as its parent.  This is not a policy
    decision made by the agent at $p$ — it is a structural constraint enforced
    by the \Fact{} layer pre-commit hook before the child node is provisioned.
    The consequence is that \textit{every node in the spawn tree, regardless of
    depth, points to the same human principal} as its ultimate accountability
    terminus.

  \item[(b)] \textbf{Escalation terminus.}  The Bailout Protocol traverses the
    parent-pointer chain $\alpha(n), \alpha(\alpha(n)), \ldots$ until reaching
    $h(n)$.  Because $h(n)$ is always a human principal and never a machine
    actor, the escalation path never terminates at an intermediate machine
    node.  The \Purpose{} layer at each node is the last machine layer before
    the human anchor — it holds the principal's intent, not a delegated subset
    of it.  This structural property makes the upstream accountability guarantee
    non-vacuous: the terminus is defined at the \Purpose{} layer, which cannot
    be replicated or substituted by any machine actor.

  \item[(c)] \textbf{Scope adjudication.}  When the \Bounds{} layer detects a
    delegation scope inheritance violation, the exception is escalated to the
    \Purpose{} layer of the parent node for adjudication.  The \Purpose{} layer
    alone has authority to redefine $R(p)$ or to authorize an exception to the
    inheritance constraint on behalf of the human principal.  No machine actor
    may grant this authorization.
\end{enumerate}

The combined effect of (a)–(c) is that the \Purpose{} layer serves as a double
anchor: it anchors the spawn tree's accountability chain at the top (human
principal $\rightarrow$ root node's \Purpose{} layer), and it anchors every
escalation path at its terminus ($h(n)$, reached through the parent-pointer
chain).  The spawn tree is not a tree of autonomous agents — it is a tree of
accountability structures, each anchored to the same human principal.

% ---------------------------------------------------------------------------
\subsection{Bounds Engine: Depth Enforcement}
\label{sec:rs-bounds-depth}

The \Bounds{} layer enforces the Hierarchy Finiteness Invariant as a
deterministic pre-commit check on every spawn event.  The enforcement mechanism
is the depth-evaluation function $\text{depth-ok}(c, p)$ that computes
$d(c) = d(p) + 1$ and compares the result against the deployment constant
$D_{\max}$.

When a node in state \texttt{CAN\_SPAWN} receives an $e_{\text{spawn}}$ event,
it transitions to \texttt{SPAWN\_PENDING} and issues a depth-confirmation
request to the \Bounds{} layer.  The \Bounds{} layer evaluates
$\text{depth-ok}(c, p)$; if the check passes, it emits $e_{\text{depth-ok}}$.
If the check fails, the \Bounds{} layer emits $e_{\text{depth-violation}}$,
transitions the spawning state machine to \texttt{SPAWN\_LOCKED}, triggers the
Bailout Protocol, and records the violation event in the Forensic Ledger.
The depth check is a deterministic integer comparison — no probabilistic
reasoning, no model judgment, no interpretation.  Either $d(c) \leq D_{\max}$
or it does not; the \Bounds{} layer enforces the condition without discretion.

The critical property of this integration is that $D_{\max}$ is a deployment
parameter established at system initialization and is not negotiable at runtime.
No machine actor, including the \Bounds{} layer itself, may raise $D_{\max}$
to accommodate a non-compliant spawn.  The deployment designer sets $D_{\max}$
once based on the operational requirements of the deployment context —
specifically, the maximum acceptable escalation path length for the Bailout
Protocol's liveness guarantee.  The \Bounds{} layer's enforcement ensures that
this contract between deployment designer and runtime system is never violated.

The delegation scope inheritance check is also enforced by the \Bounds{} layer
as a deterministic set comparison: $R(c) \subset R(p)$.  If the subset relation
fails, the \Bounds{} layer emits $e_{\text{scope-violation}}$, transitions the
spawning state machine to \texttt{SPAWN\_LOCKED}, and triggers the Bailout
Protocol.  The scope-inheritance check and the depth check are evaluated in
sequence — scope first, then depth — and both must pass before the spawn
proceeds to \texttt{CHILD\_PROVISIONED}.  This sequencing ensures that the first
detected violation is the one recorded in the Forensic Ledger, which matters for
accountability attribution when the human anchor reviews the failure.

% ---------------------------------------------------------------------------
\subsection{Fact Layer: Forensic Ledger}
\label{sec:rs-fact-fledger}

The \Fact{} layer provides two structural guarantees for Recursive Spawning:
immutable parent-pointer enforcement and complete, cryptographically chained
event recording.

\medskip
\noindent\textbf{Immutable parent pointer.}  The parent-pointer field
$\alpha(v)$ is a read-only property of the node's \Fact{} layer worksheet.
Any attempt to write to $\alpha(v)$ at runtime is intercepted by the pre-commit
hook and rejected.  The rejection event is recorded in the Forensic Ledger as
a \texttt{parent-pointer-write-denied} entry, and the Bailout Protocol is
triggered if the attempted write was initiated by a machine actor.  This
enforcement ensures that the accountability chain is never redirected after node
creation — the chain that exists at provisioning is the chain that persists
throughout the node's operational lifetime.

\medskip
\noindent\textbf{Cryptographically chained forensic record.}  Every spawn event,
every state transition in the spawning state machine, every depth-confirmation
result, every scope-inheritance evaluation, and every human directive issued in
response to a spawn-related exception is recorded as a Forensic Ledger entry.
Each entry $f$ carries the SHA-256 hash of the preceding entry, forming an
append-only chain from the genesis entry to the most recent.  The cryptographic
chaining ensures that no entry can be modified, deleted, or inserted without
breaking the chain — a condition detectable by any auditor without access to
the ledger's contents.

The Forensic Ledger entries relevant to Recursive Spawning include:
\begin{itemize}\itemsep=0pt
  \item \texttt{spawn-proposed}$(p, c, R(c), d(c))$ — parent $p$ has proposed
    to spawn child $c$ with operating range $R(c)$ and depth $d(c)$;
  \item \texttt{depth-confirmed}$(c, d(c))$ — the \Bounds{} layer confirmed
    $d(c) \leq D_{\max}$;
  \item \texttt{depth-rejected}$(c, d(c), D_{\max})$ — the depth check failed;
  \item \texttt{scope-confirmed}$(c, R(c) \subset R(p))$ — the delegation scope
    inheritance constraint is satisfied;
  \item \texttt{scope-rejected}$(c, R(c), R(p))$ — the scope-inheritance check
    failed;
  \item \texttt{child-provisioned}$(p, c, \alpha(c) = p, h(c) = h(p))$ — child
    created with immutable parent pointer and inherited human anchor;
  \item \texttt{parent-pointer-write-denied}$(v, \alpha_{\text{old}},
    \alpha_{\text{new}})$ — an attempt to modify $\alpha(v)$ was intercepted;
  \item \texttt{spawn-locked}$(v, \text{violation-type})$ — the spawning state
    machine entered \texttt{SPAWN\_LOCKED} due to a constraint violation.
\end{itemize}

The Forensic Ledger is not a conventional audit log that records actions
after they are taken — it is an enforcement mechanism that queries current
state before subsequent actions are permitted.  The \Fact{} layer pre-commit hook
verifies the current state of the spawning state machine before committing any
spawn event, ensuring that the ledger is the authoritative source of truth for
the runtime state of the spawn tree.

% ---------------------------------------------------------------------------
\subsection{Unified Integration Invariant}
\label{sec:rs-unified-inv}

The three-layer integration described above yields a single structural invariant
that is the foundation of all downstream correctness guarantees:

\begin{invariant}[Spawn Tree Accountability Invariant]
\label{inv:spawn-accountability}
For every node $v$ in a BX3 spawn tree, the accountability chain
$\alpha^{(k)}(v) = h(v)$ with $k \leq D_{\max}$ is:
\begin{enumerate}
  \item[(i)] \textbf{Immutable} — $\alpha(v)$ is never modified after node
    creation;
  \item[(ii)] \textbf{Complete} — every spawn event and state transition is
    recorded in the Forensic Ledger;
  \item[(iii)] \textbf{Terminating} — the chain terminates at a human
    principal $h(v)$, never at a machine actor;
  \item[(iv)] \textbf{Enforceable} — any violation of (i), (ii), or (iii)
    triggers the Bailout Protocol before the system state advances.
\end{enumerate}
\end{invariant}

This invariant is not a policy constraint or an operational guideline.  It is
a structural property enforced by the pre-commit hooks of all three layers
operating in concert.  The integration of Recursive Spawning with the three-layer
architecture is what makes the invariant hold — none of the three layers alone
is sufficient.  The \Purpose{} layer provides the human anchor; the \Bounds{}
layer provides depth and scope enforcement; the \Fact{} layer provides the
immutable parent pointer and the cryptographically chained forensic record.
Together they enforce Invariant~\ref{inv:spawn-accountability}, and it is this
invariant that enables the correctness properties established in the next
section.